\chapter{Polímeros Termoplásticos}
\label{cap:termoplasticos}

Los polímeros termoplásticos representan el mayor volumen de consumo en la industria de plásticos global debido a su versatilidad, fusibilidad y facilidad de conformado. Estos materiales se clasifican tradicionalmente en dos grandes grupos comerciales: plásticos de gran consumo (\textit{commodities}) y plásticos de ingeniería avanzados.

\section{Polímeros Termoplásticos de Gran Consumo (Commodities)}

\subsection{Poliolefinas: Polietileno (PE) y Polipropileno (PP)}
Las poliolefinas son hidrocarburos puros con excelente resistencia química intrínseca. El Polietileno comercial se produce con arquitecturas de ramificación bien definidas que modifican su grado de cristalinidad y densidad:
\begin{description}
    \item[HDPE (Polietileno de Alta Densidad):] Cadena lineal pura con muy pocas ramificaciones. Posee una alta cristalinidad ($70 - 90\%$) y excelente rigidez mecánica, idónea para envases rígidos de detergentes y tuberías estructurales.
    \item[LDPE (Polietileno de Baja Densidad):] Cadena con abundantes y largas ramificaciones complejas. Esto inhibe el plegamiento lineal cristalino, dando lugar a una cristalinidad del $40 - 60\%$, gran flexibilidad mecánica y excelente transparencia, ideal para películas de empaque de bolsas plásticas.
\end{description}

\subsection{Policloruro de Vinilo (PVC)}
El PVC contiene átomos de cloro electronegativos en su cadena, haciéndolo intrínsecamente rígido e inestable al procesado térmico puro. Sin embargo, su enorme versatilidad comercial radica en el papel de los aditivos y plastificantes, los cuales actúan reduciendo la $T_g$ original de $80^\circ\text{C}$ a temperatura ambiente o inferior, permitiendo obtener desde perfiles rígidos de ventanas hasta mangueras flexibles.

\section{Termoplásticos de Ingeniería y de Alto Rendimiento}

\subsection{Poliamidas (Nailon 6 y Nailon 6,6)}
Las poliamidas son polímeros de ingeniería tenaces con excelentes propiedades mecánicas y resistencia al desgaste debido al autoensamblaje espontáneo de puentes de hidrógeno intermoleculares fuertes entre los grupos amida contiguos de la cadena principal.

\subsection{Poliéter éter cetona (PEEK) e Imidas}
El PEEK es un termoplástico de alto rendimiento que posee anillos aromáticos muy rígidos y grupos funcionales éter y cetona. Exhibe una excelente temperatura de uso continuo en ambientes de trabajo extremos (superior a $250^\circ\text{C}$), siendo el sustituto predilecto de metales estructurales en la industria aeroespacial y en implantes de prótesis biomédicas.
